% ECONOMICS_PROBLEM: {"id": "C71-421", "title": "Structural enumeration of minimal balanced collections", "jel_code": "C71", "source_id": "OP-0228"}
% !TeX program = xelatex
\documentclass[11pt,a4paper]{article}
\usepackage[margin=23mm]{geometry}
\usepackage{fontspec}
\setmainfont{Latin Modern Roman}
\usepackage{amsmath,amssymb,mathtools}
\usepackage{xcolor,enumitem,booktabs,longtable,array,xurl,hyperref}
\definecolor{muted}{HTML}{53636C}
\hypersetup{colorlinks=true,urlcolor=blue,linkcolor=blue}
\setlength{\parindent}{0pt}
\setlength{\parskip}{5pt}
\newcommand{\R}{\mathbb R}
\newcommand{\E}{\mathbb E}
\newcommand{\N}{\mathbb N}
\newcommand{\ind}{\mathbf 1}
\newcommand{\Var}{\operatorname{Var}}
\newcommand{\Cov}{\operatorname{Cov}}
\newcommand{\supp}{\operatorname{supp}}
\DeclareMathOperator*{\argmin}{arg\,min}
\DeclareMathOperator*{\argmax}{arg\,max}
\newcommand{\entrymeta}[3]{{\sffamily\small\color{muted}\textbf{\texttt{#1}}\quad|\quad #2\par #3\par}}
\newcommand{\Problemref}[1]{\texttt{#1}}
\newcommand{\JELref}[2]{\texttt{#2} (JEL \texttt{#1})}
\begin{document}
\section*{Structural enumeration of minimal balanced collections}
\textbf{Problem ID:} \texttt{C71-421}\quad \textbf{JEL:} \texttt{C71}\par
% BEGIN ECONOMICS STATEMENT
\label{problem:OP-0228}
% GAME_THEORY_PROBLEM: {"local_id": "GT-098", "original_number": 98, "kind": "P", "entry_kind": "statement_only", "published": false, "review_required": true, "category": "game_theory"}
\label{game:GT-098}
{\sffamily\small\color{muted}
\noindent\textbf{ID:} GT-098. \textbf{Category:} Cooperative games and enumerative combinatorics.
\textbf{Status:} P; structural enumeration direction in the inspected source.
\par}
\subsection*{Definitions and mathematical statement}
Let $N=\{1,\ldots,n\}$, $n\geq1$. A collection is a nonempty family
$\mathcal B\subseteq2^N\setminus\{\varnothing\}$. It is \emph{balanced} if there are strictly positive real weights $(\lambda_S)_{S\in\mathcal B}$ such that
\[
 \sum_{S\in\mathcal B:\ i\in S}\lambda_S=1\qquad(i\in N),
 \quad\text{equivalently}\quad
 \sum_{S\in\mathcal B}\lambda_S\mathbf1_S=\mathbf1_N.
\]
It is \emph{minimal balanced} if no proper subcollection is balanced. Let
\[
 b_n=\left|\{\mathcal B\subseteq2^N\setminus\{\varnothing\}:
                         \mathcal B\text{ is minimal balanced}\}\right|.
\]
Find a structural counting formula or recurrence for $b_n$, including a combinatorial description explaining how minimal collections on different player-set sizes are related. Here collections are labelled, unordered families without repeated coalitions, and weights are witnesses rather than separately counted objects.

\subsection*{Short English statement}
Count the irreducible positive weighted coalition covers of a labelled player set.

\subsection*{Variants and scope boundaries}
The singleton collection $\{N\}$ is included above; excluding it changes the count by one. Counting up to player permutation is different. Ordinary set partitions are examples, but do not exhaust balanced collections. The incidence-vector definition already gives a finite exhaustive enumeration: inspect coalition families and solve the corresponding linear feasibility problems. Therefore ``effective'' cannot mean merely computable; the research request is for informative structural enumeration beyond this brute-force procedure.

\subsection*{Source relationship and status}
The source defines minimal balanced collections in Section 2 and states in the introduction that their enumeration remains open. Its convention for nontrivial minimal collections excludes $\{N\}$ in the relevant set $\mathcal B^*(n)$; the count here explicitly includes it. No closed form or new recurrence is asserted.

\subsection*{References}
\begingroup\footnotesize\raggedright
\noindent [S1] Pedro Garcia-Segador, Michel Grabisch, and Pedro Miranda, \emph{On the set of balanced games} (2025), arXiv:2501.14341v1. Inspected locators: Introduction, enumeration discussion; Section 2, balanced and minimal balanced collections and the convention for $\mathcal B^*(n)$.
\url{https://arxiv.org/pdf/2501.14341v1}
\par\endgroup

\subsection*{Common conventions: Source mathematical conventions}

\textbf{Finite games.} For a finite set $A$, $\Delta(A)$ is its probability
simplex. Players choose independent mixed strategies in Nash equilibrium;
correlation is allowed only when explicitly part of the solution concept.
In a game with payoff functions $u_i$ and strategy profile $x$, an additive
$\varepsilon$ Nash equilibrium satisfies
\[
 u_i(x)\geq u_i(x_i',x_{-i})-\varepsilon
 \quad\text{for every player $i$ and every admissible deviation $x_i'$}.
\]
For numerical approximation, payoffs are normalized as locally specified.
An $\varepsilon$ payoff residual is distinct from distance at most
$\varepsilon$ to the set of exact equilibria. Point convergence, convergence
of empirical play, and convergence in distance to an equilibrium set are
also distinct.

\textbf{Computational representations.} Unless stated otherwise, a complexity
question takes rational data with binary encoding; polynomial time is measured
in total bit length. A fixed-accuracy polynomial algorithm is not necessarily
polynomial in $1/\varepsilon$, and neither is necessarily polynomial in
$\log(1/\varepsilon)$. Succinct games and value, demand, or sampling oracles
must declare their interfaces. Exact real solutions require an explicit
representation; an unspecified list of arbitrary real numbers is not a
finite computational output. A claimed algorithmic barrier is meaningful
only for its specified model and reductions.

\textbf{Dynamic games.} Histories, information, observation, and allowable
randomization are part of the model. A stationary strategy depends only on
the current state; a positional strategy in a deterministic graph game
chooses one action at each controlled vertex. Nash equilibrium at an initial
state and equilibrium after every history are different requirements.
An expectation of a pathwise $\liminf$ payoff, a $\liminf$ of expected averages,
a limiting expected average, and a uniform finite-horizon equilibrium are
different criteria. None is silently substituted for another.

\textbf{Allocation and mechanisms.} A complete allocation assigns every item
exactly once unless otherwise stated. Zero-valued goods, negative values,
capacity constraints, feasibility, lotteries, and copies of goods affect
fairness definitions and are specified locally. Dominant-strategy incentive
compatibility, Bayesian incentive compatibility, universal truthfulness,
and truthfulness in expectation impose different inequalities. Whenever
an objective ratio has zero denominator, its convention is stated or those
instances are excluded.

\textbf{Infinite-dimensional models.} Expectations, integrals, stochastic
equations, and optimization objectives require their local regularity,
integrability, filtrations, and admissible domains. A backward--forward
mean-field system is a boundary-value problem; it is not an initial-value
semiflow without an additional construction. Long-time, large-population,
and vanishing-noise limits require an order of limits and a convergence
topology. If a minimum need not be attained, the objective is its infimum
or supremum and the approximation requirement is stated separately.
% END ECONOMICS STATEMENT
\end{document}
